%%=============================================================
%% Capítulo 5 — Desarrollo y resultados
%%=============================================================
\chapter{Desarrollo y resultados}
\label{chap:desarrollo}

Este capítulo describe en detalle el proceso de desarrollo del sistema de análisis
automático de muestras de lenguaje asistido, siguiendo la arquitectura de
\emph{pipeline} definida en el capítulo anterior. Para cada módulo del sistema se
presenta el problema que resuelve, las decisiones de diseño e implementación adoptadas
y los resultados obtenidos sobre el corpus de evaluación. Al final del capítulo se
recogen las limitaciones identificadas durante el desarrollo y las líneas de trabajo
futuro que permitirían superarlas.


%%=============================================================
\section{Visión global del sistema}
\label{sec:vision_global}
%%=============================================================

El sistema desarrollado en este trabajo transforma la grabación en vídeo de una sesión
de comunicación asistida en un formulario \ac{PAALSS} parcialmente cumplimentado. Para
ello, se articula un \emph{pipeline} de procesamiento secuencial compuesto por cuatro
etapas funcionales que se ejecutan de forma encadenada: la extracción automática de
turnos comunicativos a partir del canal de audio, el análisis lingüístico de esos
turnos según las cuatro dimensiones del protocolo, la generación de métricas agregadas
y, por último, el pre-rellenado del formulario mediante un módulo de generación
aumentada por recuperación (\ac{RAG}).

La Figura~\ref{fig:arquitectura} muestra la arquitectura completa del sistema. La
columna izquierda del diagrama representa el canal de audio, que es el único canal
implementado en este prototipo. El canal de vídeo, destinado a recuperar los
pictogramas seleccionados en el comunicador, forma parte del diseño del sistema pero
queda fuera del alcance de este trabajo por las razones expuestas en el capítulo de
objetivos. La columna central representa la capa de fusión, que integra las señales de
los distintos módulos en una representación unificada de Turno Comunicativo (\ac{CT}).
La columna derecha corresponde a las etapas de análisis lingüístico y generación de
salida.

\begin{figure}[h]
  \centering
  \missingfigure{Arquitectura general del \emph{pipeline} de análisis automático}
  \caption[Arquitectura del \emph{pipeline} PAALSS]{Arquitectura general del
    \emph{pipeline} de análisis automático. El canal de audio (izquierda) alimenta
    la capa de fusión de \ac{CT} (centro), que a su vez alimenta el análisis
    lingüístico y el módulo \ac{RAG} (derecha). El canal de vídeo queda fuera del
    alcance de este prototipo.}
  \label{fig:arquitectura}
\end{figure}

El contrato de datos entre etapas se materializa en un esquema \ac{JSON} estándar
denominado \texttt{*\_cts\_paalss.json}, definido mediante modelos \emph{Pydantic v2}
en el módulo \texttt{src/cts\_schema.py}. Este esquema recoge, para cada \ac{CT} de la
aprendiz, sus marcas temporales de inicio y fin, la fuente de recuperación, la
modalidad comunicativa, el texto transcrito y, una vez completada la etapa de análisis,
la estructura \texttt{LinguisticAnalysis} con los resultados de las cuatro dimensiones
\ac{PAALSS}. La definición explícita de este contrato garantiza que cada etapa pueda
desarrollarse y evaluarse de forma independiente sin romper la integración con el resto
del sistema.

La Tabla~\ref{tab:etapas_pipeline} resume las etapas del \emph{pipeline}, los módulos
que las implementan y los objetivos específicos a los que contribuyen.

\begin{table}[h]
  \footnotesize
  \renewcommand{\arraystretch}{1.3}
  \setlength{\tabcolsep}{4pt}
  \begin{center}
  \begin{tabular}{|L{0.5cm}|L{3.5cm}|L{4.5cm}|L{2.5cm}|}
  \hline
  \rowcolor{tema!10}
  \bft{\#} & \bft{Etapa} & \bft{Módulo} & \bft{OE asociado} \\ \hline
  1 & Transcripción \ac{ASR}        & \texttt{faster-whisper large-v3}          & OE1 \\ \hline
  2 & Diarización de hablantes      & \texttt{pyannote.audio} + DiCoW           & OE1 \\ \hline
  3 & Recuperación multicanal       & Eco + PANNs (\texttt{audio\_events.py})   & OE1 \\ \hline
  4 & Análisis lingüístico          & \texttt{spaCy es\_core\_news\_lg}         & OE2, OE3 \\ \hline
  5 & Pre-rellenado del protocolo   & \ac{RAG} + \emph{FAISS} + \emph{Ollama}  & OE4 \\ \hline
  \end{tabular}
  \end{center}
  \caption[Etapas del \emph{pipeline}]{Resumen de las etapas del \emph{pipeline}, los
    módulos que las implementan y los objetivos específicos a los que contribuyen.}
  \label{tab:etapas_pipeline}
\end{table}


%%=============================================================
\section{Etapa 1: Transcripción automática del habla}
\label{sec:etapa_asr}
%%=============================================================

\subsection{Descripción y motivación}
\label{subsec:asr_desc}

La primera etapa del \emph{pipeline} consiste en obtener una transcripción completa del
audio de la sesión con marcas temporales precisas para cada segmento de habla. Esta
transcripción es la materia prima sobre la que se construyen todas las etapas
posteriores: sin ella no es posible atribuir palabras a hablantes, detectar los
enunciados del comunicador ni calcular las métricas lingüísticas del protocolo
\ac{PAALSS}. La necesidad de procesamiento completamente local, impuesta por la
restricción \ac{RGPD} sobre datos de menores, descarta el uso de servicios de
reconocimiento de voz en la nube y obliga a emplear modelos que se ejecuten en el
hardware del clínico.

Para esta tarea se selecciona \emph{Whisper} en su variante \texttt{large-v3}, el
modelo de reconocimiento automático del habla (\ac{ASR}) publicado por Radford et al.
en 2023~\cite{whisper_radford2023}. \emph{Whisper} es un modelo \emph{transformer}
entrenado sobre 680\,000 horas de audio multilingüe procedente de Internet, con
especial cobertura del español. La ejecución local se realiza mediante la biblioteca
\texttt{faster-whisper}~\cite{whisper_radford2023}, que aplica cuantización en punto
flotante de 16 bits (\texttt{float16}) para reducir el consumo de memoria de \ac{GPU}
sin pérdida apreciable de exactitud. Con la tarjeta disponible (\emph{NVIDIA RTX 3060},
6~GB de VRAM), una sesión de 28 minutos se transcribe en aproximadamente 4--6 minutos.


\subsection{Implementación}
\label{subsec:asr_impl}

El proceso de transcripción comienza con la extracción del canal de audio del vídeo
mediante \texttt{ffmpeg}, que convierte el audio a formato WAV mono de 16~kHz, que es
la frecuencia de muestreo esperada por \emph{Whisper}. A continuación, el modelo
transcribe el audio completo de la sesión y produce una lista de segmentos, cada uno
con su texto, su marca temporal de inicio y su marca temporal de fin.

Un elemento clave de esta implementación es la construcción dinámica de un
\emph{prompt} inicial (\emph{initial\_prompt}) que se suministra al modelo antes de la
transcripción. En las sesiones de comunicación asistida con \ac{SAAC}, el vocabulario
empleado por la aprendiz está acotado al conjunto de pictogramas disponibles en su
tablero de comunicación \ac{ARASAAC}. Al incluir en el \emph{prompt} los términos más
frecuentes de ese vocabulario, el modelo reduce significativamente los errores
fonéticos sobre palabras propias del lenguaje asistido, como transcribir ``yoguri'' en
lugar de ``yogurt'' o ``era su nave'' en lugar de ``érase una vez''. El módulo
\texttt{src/vocabulary.py} gestiona la construcción de este \emph{prompt} a partir de
un fichero de vocabulario \ac{ARASAAC} de 13\,133 pictogramas, extraído mediante la
\ac{API} pública del catálogo aragonés en junio de 2026.

Adicionalmente, el mismo módulo implementa una función de corrección fonética
(\texttt{correct\_transcript}) basada en distancia de edición difusa (\emph{Levenshtein})
mediante la biblioteca \texttt{rapidfuzz}. Esta función detecta posibles errores de
\emph{Whisper} en el texto transcrito y los sustituye por el término correcto del
vocabulario \ac{ARASAAC} más próximo, siempre que la distancia de edición sea inferior
a un umbral proporcional a la longitud de la frase. El umbral proporcional es
necesario para evitar que frases muy cortas como ``Sí.'' o ``Ya.'' se colapsen
erróneamente con frases formales del tablero de comunicación.

El resultado de esta etapa es un fichero \ac{JSON} almacenado en
\texttt{results/transcription/} que contiene la lista completa de segmentos
transcritos con sus marcas temporales. Este fichero es consumido por la etapa de
diarización para la recuperación de los enunciados del comunicador mediante eco
(Sección~\ref{sec:etapa_recuperacion}).


\subsection{Resultados}
\label{subsec:asr_resultados}

La Tabla~\ref{tab:asr_ejemplos} muestra una selección de segmentos transcritos por
\emph{Whisper} en la sesión 1 («¿Dónde están las hadas y la luna?»), junto con la
corrección aplicada por el módulo de vocabulario cuando procede. Los segmentos
corresponden tanto a la voz de la interlocutora (Sol Solís) como a la voz sintetizada
del comunicador de Lucía.

\begin{table}[h]
  \footnotesize
  \renewcommand{\arraystretch}{1.3}
  \setlength{\tabcolsep}{4pt}
  \begin{center}
  \begin{tabular}{|L{1.2cm}|L{1.2cm}|L{5.5cm}|L{4.0cm}|}
  \hline
  \rowcolor{tema!10}
  \bft{Inicio (s)} & \bft{Fin (s)} & \bft{Transcripción Whisper} & \bft{Corrección} \\ \hline
    0.0 &  21.4 & Ya estamos con el cuento de los cuentos favoritos... & --- \\ \hline
   58.0 &  59.5 & Trabajar. & --- \\ \hline
  112.7 & 115.4 & Colorín colorado, colorín colorado. & --- \\ \hline
  122.1 & 126.9 & Con colorín colorado, eras una vez. & érase una vez \\ \hline
  438.4 & 438.8 & Pim Pim Pim & --- \\ \hline
  522.3 & 532.6 & (sin voz / evento acústico) & --- \\ \hline
  \end{tabular}
  \end{center}
  \caption[Ejemplos de transcripción Whisper]{Selección de segmentos transcritos en la
    sesión 1. La columna ``Corrección'' indica las sustituciones aplicadas por el
    módulo de vocabulario \ac{ARASAAC}.}
  \label{tab:asr_ejemplos}
\end{table}

El tiempo total de procesamiento de la sesión 1 (28 minutos de audio) con el modelo
\texttt{large-v3} en la \ac{GPU} disponible es de aproximadamente 5 minutos, lo que
supone una reducción del orden de 5--6$\times$ respecto al tiempo real de la sesión y
sitúa la transcripción automática muy por debajo del coste manual, que según el protocolo
\ac{PAALSS}~\cite{soto2020protocolo} puede oscilar entre 3 y 8 horas para una sesión
de 30 minutos.


%%=============================================================
\section{Etapa 2: Diarización de hablantes y segmentación en turnos}
\label{sec:etapa_diarizacion}
%%=============================================================

\subsection{Descripción y motivación}
\label{subsec:diar_desc}

La transcripción producida por \emph{Whisper} no distingue quién habla en cada
momento: el texto de la interlocutora y el de la aprendiz aparecen mezclados en la
misma secuencia de segmentos. La diarización de hablantes es el proceso de asignar
automáticamente cada segmento de audio a su hablante correspondiente, condición
indispensable para identificar exclusivamente los turnos comunicativos de la aprendiz
que establece el protocolo \ac{PAALSS}.

El escenario de diarización en las sesiones de comunicación asistida presenta dos
particularidades que lo distinguen de los escenarios para los que se han entrenado los
modelos estándar. En primer lugar, el número de hablantes es siempre dos: la aprendiz
y el interlocutor de referencia. En segundo lugar, los turnos de la aprendiz son
extremadamente cortos (duración media inferior a 2 segundos) y a menudo contienen
únicamente eventos acústicos no verbales (vocalizaciones, risas, balbuceos), lo que
dificulta la separación de clústeres basada en características acústicas de la voz.


\subsection{Implementación}
\label{subsec:diar_impl}

Para la diarización se emplea el modelo \emph{pyannote.audio} en su configuración de
\emph{pipeline} preentrenada \texttt{speaker-diarization-3.1}~\cite{pyannote_plaquet2023},
que combina detección de actividad de voz (\ac{VAD}), segmentación de hablantes y
\emph{clustering} espectral. Dado que el número de hablantes es conocido de antemano
(dos: Lucía y la interlocutora), se fija el parámetro \texttt{num\_speakers=2}, lo
que evita errores de estimación automática del número de participantes.

La fusión de la diarización con la transcripción de \emph{Whisper} se implementa
siguiendo el enfoque DiCoW (\emph{Diarization-Conditioned
Whisper})~\cite{dicow_diarization}: en lugar de alinear la transcripción completa con
la diarización mediante coincidencia temporal de segmentos, se transcribe el audio de
cada hablante de forma independiente, pasando a \emph{Whisper} únicamente los
fragmentos de audio asignados a ese hablante por el diarizador. Esta estrategia mejora
significativamente la exactitud de la transcripción sobre los segmentos cortos de la
aprendiz, porque el modelo puede aprovechar el contexto lingüístico homogéneo de un
único hablante en lugar de intentar transcribir una mezcla.

El resultado de esta etapa se almacena en un fichero \ac{JSON} intermedio que contiene,
para cada segmento, el hablante asignado, las marcas temporales, el texto transcrito y
un indicador de si el segmento se solapa temporalmente con otro hablante. Estos
segmentos se agrupan posteriormente por hablante y por proximidad temporal para
construir los Turnos Comunicativos del documento de sesión.

La definición formal de un \ac{CT} y su esquema de validación se encuentran en el módulo
\texttt{src/cts\_schema.py}. Cada \ac{CT} es una instancia del modelo Pydantic
\texttt{CT}, que valida automáticamente restricciones como que el tiempo de fin sea
posterior al de inicio, que la confianza de transcripción esté en el rango $[0,1]$ o
que la modalidad sea un valor del enumerado \texttt{Modality} (\emph{comunicador},
\emph{vocalización}, \emph{no\_verbal}). El tipo de objeto \texttt{CTDocument}
encapsula la sesión completa y expone métodos de serialización (\texttt{save},
\texttt{load}) y de filtrado (\texttt{learner\_cts(valid\_only=True)}) que las etapas
posteriores consumen de forma uniforme.


\subsection{Filtrado y limpieza de turnos}
\label{subsec:diar_filtrado}

Un problema habitual en la diarización de sesiones naturales es la aparición de
segmentos de Lucía que, en realidad, corresponden a fragmentos de habla del
interlocutor asignados erróneamente al clúster de la aprendiz. Para mitigar este
efecto, el módulo \texttt{build\_document} de \texttt{cts\_schema.py} aplica
automáticamente una función de descarte (\texttt{discard\_reason}) sobre cada \ac{CT}
de la aprendiz. Esta función marca un turno como descartado por el motivo
\texttt{posible\_interlocutor} si contiene tres o más palabras pertenecientes al
conjunto de marcadores discursivos del habla conectada adulta (como \emph{pues},
\emph{venga}, \emph{vamos}, \emph{claro} o \emph{vale}), palabras que el comunicador
electrónico de Lucía nunca produciría. De forma análoga, los segmentos de diarización
que no contienen transcripción de \emph{Whisper} se marcan como \texttt{vacio}.

El criterio de descarte es deliberadamente conservador: se prioriza no eliminar
enunciados telegráficos reales de Lucía (un \ac{CT} válido de la aprendiz puede
contener, por ejemplo, únicamente la palabra ``vamos'') sobre eliminar todos los
posibles errores. Los \ac{CT} descartados no se eliminan del documento, sino que
permanecen en él con la marca correspondiente para que el especialista pueda
revisarlos.

La Tabla~\ref{tab:cts_sesion1} resume los resultados de esta etapa sobre la sesión 1.

\begin{table}[h]
  \footnotesize
  \renewcommand{\arraystretch}{1.3}
  \setlength{\tabcolsep}{4pt}
  \begin{center}
  \begin{tabular}{|L{7cm}|L{2.5cm}|}
  \hline
  \rowcolor{tema!10}
  \bft{Métrica} & \bft{Valor} \\ \hline
  \ac{CT} totales detectados (Lucía + interlocutor)  & 202 \\ \hline
  \ac{CT} de Lucía (diarización)                     & 54  \\ \hline
  \ac{CT} de Lucía descartados (vacíos)              & 20  \\ \hline
  \ac{CT} de Lucía descartados (posible interlocutor)& 2   \\ \hline
  \ac{CT} de Lucía válidos (canal diarización)       & 32  \\ \hline
  Mínimo \ac{PAALSS} (50 \ac{CT})                    & No alcanzado \\ \hline
  \end{tabular}
  \end{center}
  \caption[\ac{CT} de diarización sesión 1]{Resultados de la etapa de diarización y
    filtrado sobre la sesión 1. El umbral mínimo de 50~\ac{CT} válidos que establece
    el protocolo \ac{PAALSS}~\cite{soto2020protocolo} no se alcanza únicamente con el
    canal de audio, lo que motiva los módulos de recuperación complementarios de la
    Sección~\ref{sec:etapa_recuperacion}.}
  \label{tab:cts_sesion1}
\end{table}


%%=============================================================
\section{Etapa 3: Recuperación multicanal de turnos comunicativos}
\label{sec:etapa_recuperacion}
%%=============================================================

\subsection{Descripción y motivación}
\label{subsec:rec_desc}

Los resultados de la etapa anterior revelan un problema estructural del canal de audio
en el contexto de la comunicación asistida: el diarizador únicamente recupera 32~\ac{CT}
válidos de Lucía, cuando el estándar de referencia (\texttt{PAALSS2\_lucia.pdf})
contiene 44 \ac{CT} anotados por la especialista. Un análisis de la brecha entre ambas
cifras permite identificar dos categorías de turnos perdidos.

La primera categoría corresponde a los enunciados producidos mediante el comunicador
electrónico de Lucía. El dispositivo emite una voz sintetizada adulta y fluida que el
diarizador \emph{pyannote} asigna sistemáticamente al clúster del interlocutor, dado
que acústicamente es indistinguible de la voz de la profesional de referencia. Estos
enunciados contienen la mayor parte del vocabulario de la aprendiz y su pérdida limita
gravemente la utilidad del análisis lingüístico posterior.

La segunda categoría corresponde a vocalizaciones no verbales de la aprendiz (risas,
balbuceos, sonidos expresivos) que el diarizador detecta como segmentos sin
transcripción, y que el paso de filtrado descarta como \texttt{vacio}. Sin embargo,
estos eventos son comunicativamente relevantes y deben recogerse en el protocolo
\ac{PAALSS} como anotaciones de la columna de eventos no verbales.

Para mitigar ambas carencias se han diseñado dos módulos de recuperación complementarios
que operan sobre el resultado de la diarización.


\subsection{Módulo de recuperación por eco (canal comunicador)}
\label{subsec:rec_eco}

El módulo de recuperación por eco (\emph{echo detection}) explota la práctica clínica
habitual de que el interlocutor repite en voz alta el mensaje que la aprendiz ha
seleccionado en el comunicador, inmediatamente después de que el dispositivo lo emita.
Esta repetición actúa como señal de confirmación del enunciado y es audible en el
canal de audio aunque la voz sintetizada original haya sido incorrectamente asignada al
clúster del interlocutor.

El algoritmo opera en tres pasos. En primer lugar, se recorre la transcripción bruta
producida por \emph{Whisper} en la etapa 1 y se puntúan todos los segmentos según la
función \texttt{score\_as\_device\_tts} del módulo \texttt{src/vocabulary.py}. Esta
función calcula una puntuación compuesta entre 0 y 1 que combina cuatro indicadores:
la duración del segmento (los enunciados del comunicador suelen ser breves), el número
de palabras, la proporción de palabras presentes en el vocabulario \ac{ARASAAC} y la
coincidencia con frases formales del tablero de comunicación. Los segmentos que superan
un umbral de puntuación de 0{,}45 se consideran candidatos a enunciado del comunicador.

En segundo lugar, para cada candidato se busca en la transcripción diarizada del
interlocutor, en una ventana de hasta 8 segundos posteriores al final del candidato,
si aparece algún fragmento que contenga el mismo término en forma de cita. La aparición
de esta repetición (\emph{eco}) confirma que el candidato es un enunciado real del
comunicador y no un fragmento de habla del interlocutor transcrito erróneamente.

En tercer lugar, los candidatos confirmados por eco se añaden al documento de sesión
como \ac{CT} con la fuente \texttt{echo\_diarized}, diferenciándolos así de los
\ac{CT} obtenidos por diarización directa (\texttt{diarization}). Esta distinción de
fuente queda registrada en el campo \texttt{procedencia} del esquema \ac{JSON} y
permite al especialista revisar con mayor detención los \ac{CT} recuperados por este
método.

El criterio de confirmación es deliberadamente conservador, exigiendo puntuación alta
y eco posterior, para evitar incorporar falsos positivos al análisis lingüístico. Una
variante más permisiva que incluía eco bidireccional y mayor tolerancia en la distancia
de edición fue probada y descartada porque incrementaba el \emph{recall} en
aproximadamente un 10~\% pero introducía unos 7 falsos positivos por sesión (términos
como \emph{casa}, \emph{río} o \emph{ver} que correspondían a habla del interlocutor,
no a enunciados del comunicador).

Sobre la sesión 1, el módulo de eco conservador recupera 3 \ac{CT} adicionales del
comunicador (\emph{Aquí}, \emph{Mira} y \emph{Si no están}), lo que eleva el total de
\ac{CT} válidos de Lucía de 32 a 35.


\subsection{Módulo de eventos no verbales (PANNs)}
\label{subsec:rec_panns}

Los eventos acústicos no verbales de Lucía se detectan y clasifican mediante el módulo
\texttt{src/audio\_events.py}. Los candidatos a evento no verbal se identifican como
los segmentos asignados a Lucía por \emph{pyannote} que tienen una duración mínima de
0{,}4 segundos y no coinciden temporalmente con ninguna palabra transcrita por
\emph{Whisper}. El uso de \emph{pyannote} como detector de actividad de voz (\ac{VAD})
es suficiente para capturar vocalizaciones expresivas (risas, balbuceos), dado que
\emph{pyannote} está entrenado para detectar habla humana en sentido amplio y captura
estas emisiones sin necesidad de un \ac{VAD} complementario basado en energía de señal.

La clasificación de cada evento utiliza el modelo \textbf{PANNs CNN14} preentrenado en
AudioSet, un modelo de reconocimiento de audio de propósito general con 527 clases, que
se ejecuta completamente de forma local sobre PyTorch. Para cada candidato, el módulo
extrae el clip de audio correspondiente, lo remuestrea a 32~kHz (la frecuencia esperada
por PANNs) y obtiene un vector de probabilidades sobre las 527 clases. A continuación,
aplica un mapeo curado (\texttt{AUDIOSET\_TO\_PAALSS}) que traduce las etiquetas de
AudioSet a las categorías de eventos no verbales empleadas en el protocolo \ac{PAALSS}:
\emph{risa}, \emph{llanto}, \emph{balbuceo}, \emph{suspiro}, \emph{grito} y
\emph{queja}. Los eventos que no superan el umbral de confianza de 0{,}5 se marcan
con \texttt{revisar=True} para que el especialista los verifique manualmente, ya que
AudioSet es un corpus de dominio general y la voz disártrica infantil suele clasificarse
como \emph{Speech}, reduciendo la confianza de las etiquetas específicas.

Los eventos no verbales se añaden al documento de sesión como \ac{CT} con la fuente
\texttt{acoustic\_event} y modalidad \texttt{no\_verbal}. Su texto se registra como
\texttt{(no verbal: risa)} o similar, y el campo \texttt{anotacion} del esquema
recoge la etiqueta de tipo de sonido obtenida por PANNs. Esta etiqueta es una
\emph{pre-anotación} que el especialista debe confirmar o corregir en su revisión, en
línea con la filosofía del protocolo \ac{PAALSS} de mantener al especialista como
decisor final.


\subsection{Resultados de recuperación multicanal}
\label{subsec:rec_resultados}

La Tabla~\ref{tab:recall_canales} muestra el \emph{recall} acumulado de cada canal de
recuperación sobre la sesión 1, calculado frente a los 44 \ac{CT} del estándar de
referencia anotado por la especialista.

\begin{table}[h]
  \footnotesize
  \renewcommand{\arraystretch}{1.3}
  \setlength{\tabcolsep}{4pt}
  \begin{center}
  \begin{tabular}{|L{5.5cm}|L{2cm}|L{2cm}|L{3cm}|}
  \hline
  \rowcolor{tema!10}
  \bft{Canal} & \bft{\ac{CT} recuperados} & \bft{\emph{Recall}} & \bft{Fuente JSON} \\ \hline
  Diarización (\emph{pyannote} + Whisper)  & 32 & 18\,\% & \texttt{diarization} \\ \hline
  + Eco del comunicador                    & 35 & $\approx$\,36\,\%  & \texttt{echo\_diarized} \\ \hline
  + Eventos no verbales (PANNs)            & 38 & $\approx$\,40\,\%  & \texttt{acoustic\_event} \\ \hline
  \bft{Techo teórico canal audio}          & \bft{---} & \bft{$\approx$\,47\,\%} & --- \\ \hline
  \end{tabular}
  \end{center}
  \caption[\emph{Recall} por canal de recuperación]{Recall acumulado de los tres canales
    de recuperación sobre la sesión 1 («¿Dónde están las hadas y la luna?», 44~\ac{CT}
    en el estándar de referencia). El techo teórico del canal de audio se estima como
    el porcentaje de \ac{CT} de la aprendiz cuyo contenido está presente en la
    transcripción bruta de \emph{Whisper}, incluyendo aquellos atribuidos al interlocutor.}
  \label{tab:recall_canales}
\end{table}

Estos resultados muestran que el canal de audio, incluso combinando los tres módulos de
recuperación, alcanza un techo de aproximadamente el 47~\% de los \ac{CT} del estándar
de referencia. La causa principal de este techo es estructural: aproximadamente el
56~\% de los turnos de Lucía corresponden a eventos no verbales o gestos que no dejan
huella en el canal de audio, según el análisis de la distribución de turnos de la
sesión. Los turnos que escapan al canal de audio son aquellos producidos exclusivamente
mediante selección de pictogramas en el comunicador sin activación de la síntesis de
voz, mediante gestos oculares o manuales, o mediante señalización directa. Estos
enunciados solo son recuperables mediante el procesamiento del canal de vídeo, que
constituye la extensión más prioritaria del sistema descrita en la
Sección~\ref{sec:limitaciones}.


%%=============================================================
\section{Etapa 4: Análisis lingüístico PAALSS}
\label{sec:etapa_nlp}
%%=============================================================

\subsection{Descripción y motivación}
\label{subsec:nlp_desc}

Una vez construido el documento de sesión con los \ac{CT} de la aprendiz, la cuarta
etapa aplica el análisis lingüístico que corresponde a las cuatro dimensiones del
protocolo \ac{PAALSS}: vocabulario, morfología, longitud media del enunciado y
sintaxis. Esta etapa opera exclusivamente sobre los \ac{CT} válidos de la aprendiz
(campo \texttt{descartado=False}) y enriquece cada uno de ellos con una estructura
\texttt{LinguisticAnalysis} que queda almacenada en el campo \texttt{analisis} del
esquema \ac{JSON}.

El análisis lingüístico de lenguaje asistido plantea retos específicos que lo
diferencian del procesamiento de texto convencional. Los pictogramas del comunicador
representan formas de cita (infinitivo, singular masculino), por lo que la
ausencia de flexión morfológica no es un error sino una característica intrínseca de
la modalidad. Del mismo modo, los enunciados de la aprendiz son telegráficos y a menudo
contienen una o dos palabras, lo que hace que los análisis de dependencias sintácticas
de los modelos estándar tengan una precisión reducida sobre este tipo de texto.


\subsection{Implementación}
\label{subsec:nlp_impl}

El módulo de análisis lingüístico (\texttt{src/pipeline/etapa3\_nlp.py}) carga el
modelo \texttt{es\_core\_news\_lg} de \emph{spaCy}~\cite{spoken_spanish_pos} y lo
aplica sobre el texto de cada \ac{CT} válido. Por razones de compatibilidad con el
entorno de ejecución en Windows, el módulo se ejecuta en un subproceso Python aislado
(\texttt{src/pipeline/\_nlp\_worker.py}), separado del proceso principal que ha
cargado los modelos de \ac{GPU} de las etapas anteriores. Esta separación evita un
conflicto de runtimes \ac{OMP} entre \emph{PyTorch} y \emph{spaCy} que provoca una
terminación abrupta del proceso en Windows.

Para cada \ac{CT}, la función \texttt{\_analyze} realiza las siguientes operaciones.
En primer lugar, tokeniza el texto y filtra signos de puntuación y espacios, obteniendo
la lista de \emph{tokens} de contenido. En segundo lugar, para cada \emph{token} de
contenido extrae su etiqueta de categoría gramatical (\ac{POS}) mediante el clasificador
de \emph{spaCy} y la traduce al conjunto de categorías del protocolo \ac{PAALSS}:
\emph{verbo}, \emph{nombre}, \emph{adjetivo}, \emph{artículo}, \emph{pronombre},
\emph{preposición}, \emph{conjunción}, \emph{adverbio}, \emph{interjección} y
\emph{numeral}. En tercer lugar, calcula la \ac{LME} contando morfemas según las reglas
del protocolo: cada \emph{token} de contenido aporta un morfema base, más uno adicional
si está en plural, más uno si es un verbo conjugado en pasado o futuro, más otro si
tiene persona y número explícitos. En cuarto lugar, detecta la presencia del patrón
\ac{SVO} mediante heurísticos sobre el árbol de dependencias producido por \emph{spaCy},
identificando el sujeto (\emph{nsubj}), el verbo raíz (\emph{ROOT}) y el objeto
(\emph{obj}, \emph{iobj} u \emph{obl}) y calculando el orden relativo de sus posiciones
en el enunciado.

El resultado de este análisis para cada \ac{CT} se valida automáticamente como
instancia del modelo Pydantic \texttt{LinguisticAnalysis}, que agrupa cuatro subestructuras:
\texttt{Vocabulary} (con los campos \emph{tokens}, \emph{lemmas}, \emph{tnw} y
\emph{categories}), \texttt{MorphologicalFeature} (una lista de registros por
\emph{token}), \texttt{MLU} (con los campos \emph{morphemes} y \emph{words}) y
\texttt{Syntax} (con los campos \emph{order}, \emph{svo}, \emph{sentence\_type} y
\emph{connectors}).


\subsection{Resultados del análisis lingüístico}
\label{subsec:nlp_resultados}

La Tabla~\ref{tab:nlp_ejemplos} muestra el resultado del análisis para una selección
de \ac{CT} de Lucía en la sesión 1, incluyendo las cuatro dimensiones \ac{PAALSS}.

\begin{table}[h]
  \footnotesize
  \renewcommand{\arraystretch}{1.3}
  \setlength{\tabcolsep}{3pt}
  \begin{center}
  \begin{tabular}{|L{3.5cm}|L{1.8cm}|L{2.0cm}|L{1.0cm}|L{3.8cm}|}
  \hline
  \rowcolor{tema!10}
  \bft{Texto CT} & \bft{Categorías} & \bft{Morfología} & \bft{LME} & \bft{Sintaxis} \\ \hline
  Trabajar.
    & verbo: 1
    & infinitivo (forma base)
    & 1{,}0
    & Orden: V (solo verbo) \\ \hline
  Mira, mira, mira, mira.
    & verbo: 1
    & imperativo (correcto)
    & 3{,}0
    & Orden: V \\ \hline
  Necesito ayuda.
    & verbo: 1, nombre: 1
    & pres. 1.ª sg. + sustantivo
    & 4{,}0
    & Orden: VO \\ \hline
  No hay nada, no hay nada.
    & verbo: 1, adverbio: 2
    & negación + pres. impersonal
    & 3{,}0
    & Negativa \\ \hline
  Colorín colorado, colorín colorado.
    & nombre: 2, adjetivo: 2
    & nominales en forma base
    & 4{,}0
    & Orden: SN \\ \hline
  \end{tabular}
  \end{center}
  \caption[Análisis lingüístico de \ac{CT} seleccionados]{Resultado del análisis
    lingüístico automático sobre una selección de \ac{CT} de Lucía en la sesión 1.
    La columna LME muestra la \ac{LME} en morfemas calculada por el módulo.}
  \label{tab:nlp_ejemplos}
\end{table}

Las métricas agregadas de vocabulario para la sesión 1, calculadas sobre los 35~\ac{CT}
válidos recuperados por los tres canales, se recogen en la
Tabla~\ref{tab:metricas_vocabulario}.

\begin{table}[h]
  \footnotesize
  \renewcommand{\arraystretch}{1.3}
  \setlength{\tabcolsep}{4pt}
  \begin{center}
  \begin{tabular}{|L{7cm}|L{2cm}|}
  \hline
  \rowcolor{tema!10}
  \bft{Métrica} & \bft{Valor} \\ \hline
  TNW (Total de palabras)                    & 74   \\ \hline
  NDW (Número de palabras diferentes)        & 41   \\ \hline
  Ratio NDW/TNW                              & 0{,}55 \\ \hline
  LME media (en morfemas)                    & 2{,}8 \\ \hline
  \ac{CT} con patrón SVO completo            & 3 (8{,}6\,\%) \\ \hline
  \ac{CT} con verbo presente                 & 6 (17{,}1\,\%) \\ \hline
  \end{tabular}
  \end{center}
  \caption[Métricas PAALSS sesión 1]{Métricas lingüísticas calculadas automáticamente
    sobre los 35~\ac{CT} válidos de Lucía en la sesión 1. Los valores corresponden al
    resultado del análisis automático y no han sido validados frente al estándar de
    referencia.}
  \label{tab:metricas_vocabulario}
\end{table}

La baja frecuencia de estructuras \ac{SVO} completas (8{,}6\,\%) es coherente con la
naturaleza telegráfica del lenguaje asistido, donde los pictogramas del comunicador
representan el contenido léxico mínimo necesario para comunicar la intención y la
flexión gramatical queda sistemáticamente omitida. Este patrón es precisamente el que
el protocolo \ac{PAALSS} busca cuantificar como indicador del nivel de desarrollo
lingüístico de la aprendiz.


%%=============================================================
\section{Etapa 5: Pre-rellenado automático del formulario}
\label{sec:etapa_rag}
%%=============================================================

\subsection{Descripción y motivación}
\label{subsec:rag_desc}

La etapa final del \emph{pipeline} traduce las métricas lingüísticas calculadas y los
\ac{CT} analizados en propuestas de valores para los campos del formulario oficial
\ac{PAALSS}. Este es el producto visible del sistema para el profesional clínico: en
lugar de revisar una tabla de métricas abstractas, el especialista recibe el formulario
\ac{PAALSS} con los campos numéricos ya cumplimentados y las celdas de texto libre con
propuestas de descripción ancladas en fragmentos concretos de la sesión.

La arquitectura de generación aumentada por recuperación (\ac{RAG}) es la elección
adecuada para esta tarea porque combina las ventajas de los modelos de lenguaje grandes
(capacidad de generar texto en lenguaje natural estructurado) con la necesidad de basar
las propuestas exclusivamente en la evidencia de la sesión (los \ac{CT} analizados),
sin recurrir a conocimiento externo que podría introducir alucinaciones. Los trabajos de
Babaeipour et al.~\cite{clinical_nlp_ai_protocol} muestran que, en tareas de
extracción de información clínica a partir de protocolos estructurados, la arquitectura
\ac{RAG} mejora la exactitud del pre-rellenado de un 62{,}6~\% (modelo de lenguaje
autónomo) a un 89~\% respecto a anotaciones expertas.


\subsection{Diseño del módulo}
\label{subsec:rag_diseno}

El módulo \ac{RAG} se organiza en tres componentes. El primero es la base de datos
vectorial: los textos y análisis de los \ac{CT} de la sesión se codifican como
vectores de \emph{embeddings} mediante un modelo de codificación de oraciones en
español, y se almacenan en un índice \emph{FAISS}~\cite{ollama} que permite la
recuperación eficiente por similitud semántica. El segundo componente es el motor de
recuperación: para cada campo del formulario \ac{PAALSS} que requiere una descripción
textual, se construye una consulta semántica que recupera los \ac{CT} más relevantes de
la sesión. El tercer componente es el modelo de lenguaje local: los \ac{CT} recuperados
se concatenan como contexto y se pasan, junto con la instrucción de cumplimentar el
campo concreto, al modelo de lenguaje ejecutado mediante \emph{Ollama}~\cite{ollama}.

La instrucción al modelo incluye de forma explícita la directriz de indicar
ausencia de información cuando el contexto recuperado no contenga evidencia suficiente
para proponer un valor, de modo que el especialista pueda identificar rápidamente qué
campos requieren revisión manual adicional.

El estado de implementación actual de este módulo es de diseño validado y desarrollo
parcial. Los componentes de vectorización y recuperación \emph{FAISS} están
implementados y probados; la integración con \emph{Ollama} y la generación de las
propuestas para cada campo del formulario están pendientes, como se detalla en la
Sección~\ref{sec:limitaciones}.


%%=============================================================
\section{Interfaz de usuario}
\label{sec:interfaz}
%%=============================================================

El sistema se ofrece al usuario a través de una interfaz web local construida con
\emph{Streamlit}, que permite al profesional clínico operar el \emph{pipeline}
completo sin conocimientos de programación. La interfaz se organiza en dos zonas
principales. La zona superior permite seleccionar el vídeo de la sesión (mediante ruta
local, para evitar copiar ficheros de gran tamaño en memoria) y configurar los
parámetros básicos del análisis, como el número de hablantes o el \emph{token} de
acceso a los modelos de diarización. La zona inferior presenta los resultados una vez
completado el procesamiento, organizados en cuatro pestañas: resumen de \ac{CT},
vocabulario, morfología y análisis.

La Figura~\ref{fig:interfaz} ilustra la vista de resultados de la pestaña de resumen
de \ac{CT}, que muestra para cada turno el número de orden, el instante de inicio, la
duración, la fuente de recuperación, la modalidad y el texto. Los \ac{CT} descartados
por el filtro de limpieza se muestran en un panel colapsable separado, de modo que el
especialista puede revisarlos sin que contaminen la vista principal.

\begin{figure}[h]
  \centering
  \missingfigure{Captura de la interfaz Streamlit: pestaña de resumen de \ac{CT} de la sesión 1.}
  \caption[Interfaz Streamlit]{Vista de la pestaña de resumen de \ac{CT} en la interfaz
    \emph{Streamlit}. Cada fila muestra un turno comunicativo válido de la aprendiz con
    su fuente de recuperación y su contenido.}
  \label{fig:interfaz}
\end{figure}


%%=============================================================
\section{Limitaciones y análisis de resultados}
\label{sec:limitaciones}
%%=============================================================

\subsection{Techo del canal de audio}
\label{subsec:lim_audio}

El resultado más significativo de la evaluación sobre la sesión 1 es que el canal de
audio, incluso con los tres módulos de recuperación combinados, alcanza un \emph{recall}
de aproximadamente el 40~\% respecto al estándar de referencia de la especialista. Este
techo no es consecuencia de limitaciones evitables de los modelos empleados, sino de
una restricción estructural del dominio: una parte sustancial de los enunciados de la
aprendiz se produce exclusivamente mediante selección de pictogramas en el comunicador
sin síntesis de voz, o mediante gestos oculares, señalización o expresiones corporales,
modalidades que no dejan ninguna huella en el canal de audio.

El análisis de la distribución de turnos en la sesión 1, realizado a partir de las
anotaciones del estándar de referencia, muestra que aproximadamente el 56~\% de los
turnos comunicativos de Lucía son eventos no verbales o vocalizaciones sin transcripción
inteligible. Esto implica que el techo teórico del canal de audio es del 44~\% de los
turnos totales de la aprendiz, en perfecta concordancia con el \emph{recall} observado.
Este resultado valida empíricamente la decisión de diseño de incluir el canal de vídeo
como componente obligatorio del sistema completo (OE1), y no como mejora opcional.

\subsection{Limitaciones del módulo de recuperación por eco}
\label{subsec:lim_eco}

El módulo de recuperación por eco recupera con alta precisión los \ac{CT} del
comunicador cuya voz sintetizada es claramente distinta de la del interlocutor y va
seguida de una repetición audible. Sin embargo, existen casos en que el interlocutor
no repite el enunciado, en que la repetición ocurre fuera de la ventana de 8 segundos,
o en que el enunciado del comunicador consiste en una sola palabra que también forma
parte del habla natural del interlocutor (como \emph{sí}, \emph{no} o \emph{aquí}),
lo que hace imposible determinar la fuente sin información visual. En estos casos el
módulo no recupera el \ac{CT}, contribuyendo a la brecha entre el \emph{recall}
observado y el techo teórico.

\subsection{Limitaciones del análisis lingüístico}
\label{subsec:lim_nlp}

El análisis lingüístico basado en \emph{spaCy} está entrenado sobre texto escrito
convencional, no sobre lenguaje telegráfico ni lenguaje asistido. Los enunciados de
Lucía, compuestos con frecuencia por una o dos palabras en forma de cita, presentan
árboles de dependencias incompletos que reducen la precisión del análisis sintáctico.
La detección del patrón \ac{SVO} mediante heurísticos sobre dependencias resulta
razonablemente fiable cuando el enunciado tiene estructura predicativa completa, pero
es inaplicable sobre enunciados nominales (la mayoría). Para mejorar este componente
sería necesario aplicar modelos de clasificación de función comunicativa específicamente
entrenados sobre lenguaje asistido en español, para lo que no existe actualmente ningún
corpus etiquetado disponible en este dominio.

\subsection{Corpus de evaluación limitado}
\label{subsec:lim_corpus}

El corpus de evaluación disponible se reduce a dos sesiones anotadas manualmente por la
especialista de referencia, correspondientes a Lucía en dos actividades de lectura
compartida distintas. Aunque estas dos sesiones son suficientes para demostrar el
funcionamiento del sistema y obtener métricas iniciales de rendimiento, son insuficientes
para estimar de forma robusta la capacidad de generalización del sistema a otras
aprendices, otros interlocutores, otros tipos de actividad o diferentes configuraciones
de tablero de comunicación. La ampliación del corpus de evaluación a través de la
colaboración con la \emph{Fábrica de Palabras} es la condición más importante para
la validación clínica del sistema.

\subsection{Estado del módulo RAG}
\label{subsec:lim_rag}

El módulo de pre-rellenado del formulario mediante \ac{RAG} se encuentra en estado de
diseño validado y desarrollo parcial. Los componentes de vectorización y recuperación
están operativos; la integración con el modelo de lenguaje local y la generación de
propuestas campo a campo del formulario \ac{PAALSS} están pendientes de implementación.
En consecuencia, la evaluación cuantitativa de este módulo frente al estándar de
referencia (OE5, \emph{kappa de Cohen} campo a campo) no ha podido completarse en el
plazo de este trabajo.

\subsection{Trabajo futuro}
\label{subsec:trabajo_futuro}

Las limitaciones identificadas permiten trazar de forma concreta las líneas de trabajo
futuro con mayor impacto esperado. La extensión más prioritaria es la implementación
del canal de vídeo para la detección de pictogramas seleccionados en el comunicador,
que permitiría superar el techo del 47~\% del canal de audio y alcanzar un \emph{recall}
cercano al 90~\% de los turnos de la aprendiz. En segundo lugar, la finalización del
módulo \ac{RAG} y su evaluación cuantitativa frente a las anotaciones expertas
completaría la demostración del sistema extremo a extremo. En tercer lugar, la
ampliación del corpus a nuevas sesiones y nuevas aprendices, en colaboración con la
\emph{Fábrica de Palabras}, permitiría estimar la generalización del sistema y
orientar los esfuerzos de mejora hacia los módulos con mayor impacto clínico.


% Local Variables:
%  coding: utf-8
%  mode: latex
%  mode: flyspell
%  ispell-local-dictionary: "castellano8"
% End:
